\documentclass[11pt]{article}
\usepackage[margin=1.25in]{geometry}
\usepackage{amsmath,amssymb,amsthm}
\usepackage{hyperref}
\usepackage{booktabs}
\usepackage{graphicx}
\usepackage{xcolor}
\usepackage{microtype}
\usepackage{natbib}
\usepackage{algorithm}
\usepackage{algpseudocode}

\hypersetup{colorlinks=true, linkcolor=blue, citecolor=blue, urlcolor=blue}

\newtheorem{theorem}{Theorem}
\newtheorem{lemma}[theorem]{Lemma}
\newtheorem{corollary}[theorem]{Corollary}
\newtheorem{definition}[theorem]{Definition}
\newtheorem{remark}[theorem]{Remark}

\DeclareMathOperator{\supp}{supp}
\newcommand{\eps}{\varepsilon}
\newcommand{\C}{\mathbb{C}}
\newcommand{\R}{\mathbb{R}}
\newcommand{\N}{\mathbb{N}}
\newcommand{\E}{\mathbb{E}}
\newcommand{\Otilde}{\widetilde{O}}

\title{%
  \textbf{Adaptive Quantum Oracle Sketching with Sublinear Sample Complexity}\\[6pt]
  {\large Importance-Weighted Phase Accumulation, Hierarchical Query Reduction,}\\[2pt]
  {\large and Interferometric Classical Shadows}
}
\author{
  Tommaso R.\ Marena\\
  \textit{Department of Mathematics and Statistics}\\
  \textit{The Catholic University of America}\\
  Washington, DC 20064\\
  \texttt{marena@cua.edu}
}
\date{April 2026}

\begin{document}
\maketitle

\begin{abstract}
We present three improvements to the quantum oracle sketching framework of Zhao et\ al.\ (2025).  Their construction of a phase-oracle diagonal $U_f: |x\rangle\mapsto e^{i\pi f(x)}|x\rangle$ via expected-unitary methods requires $O(NQ^2)$ samples, where $N$ is the domain size and $Q$ is the number of linear queries.  For sparse Boolean functions with $K$-element support and for Fourier-sparse signal priors, we show this cost is not necessary.

\textbf{(1) Adaptive sparse oracle.}  A two-phase importance-sampling scheme reduces the cost to $O(KQ^2)$, an $N/K$ factor improvement.  The pilot phase estimates the support distribution in $O(K)$ samples; the main phase applies the expected-unitary construction with importance weights.  We prove a sharp concentration bound showing the error on $\supp(f)$ is $O(\sqrt{K/M_{\mathrm{main}}})$, matching the information-theoretic lower bound up to constants.

\textbf{(2) Hierarchical oracle sketching.}  By composing $k$ levels of oracle sketches, each refining the residual error of the previous, we reduce the query-dependent cost from $O(NQ^2)$ to $O(NQ^{2-1/k})$.  For any fixed $\eps > 0$ and $k = \lceil 1/\eps\rceil$, this yields $O(NQ^{2-\eps})$, strictly sub-quadratic in $Q$.

\textbf{(3) Variational warmstart.}  For signals with a Fourier-sparse prior (at most $K_F$ dominant modes), a gradient-descent pre-training step reduces the effective domain to $K_F \ll N$ Fourier modes.  Combined with (1), the total cost becomes $O(K_F Q^{2-1/k})$, achieving an $(N/K_F)\cdot Q^{1/k}$ improvement over the baseline.

All constructions are implemented and verified in open-source JAX code.  The combined bound at concrete parameters ($N = 2^{20}$, $K_F = 64$, $Q = 100$, $k = 3$) gives a $76{,}293\times$ reduction in oracle sample complexity.  We provide complete proofs, matching lower bounds for (1), and empirical validation across all four contributions.
\end{abstract}

\section{Introduction}\label{sec:intro}

The problem of constructing efficient quantum oracles from classical queries sits at the heart of quantum advantage arguments.  A phase oracle $U_f$ for a Boolean function $f:\{0,1\}^n\to\{0,1\}$ encodes $f$ as a phase: $U_f|x\rangle = e^{i\pi f(x)}|x\rangle$.  Direct implementation requires exponential circuit depth in general; approximating $U_f$ via linear measurements of $f$, however, is feasible in polynomial time under sparsity assumptions.

Zhao, Iosue, Bharti, Deshpande, and Wilde (2025) introduced a beautiful expected-unitary construction that builds an approximation to the oracle diagonal using $M = O(NQ^2)$ samples.  Their approach, which we term the \emph{uniform oracle sketch}, is information-theoretically tight when $f$ is uniformly dense.  However, for the structured functions that appear in quantum chemistry, combinatorial optimization, and signal processing---functions that are $K$-sparse, Fourier-concentrated, or hierarchically structured---the $O(N)$ dependence represents a fundamental bottleneck.

This paper addresses that bottleneck.  We ask a simple question: \emph{if we know or can estimate that $f$ has support of size $K \ll N$, can we build the oracle in $O(K)$ rather than $O(N)$ samples?}  The answer is yes, and the same intuition extends to query complexity (yielding $Q^{2-1/k}$ instead of $Q^2$) and to Fourier-sparse priors.

\subsection{Our Contributions}

Let $\eps > 0$ be a target approximation error, $N = 2^n$ the domain size, $K = |\supp(f)|$ the support size, $K_F$ the number of dominant Fourier modes, $Q$ the number of linear queries, and $k \geq 1$ the number of hierarchical levels.

\begin{center}
\begin{tabular}{lcc}
\toprule
\textbf{Method} & \textbf{Sample complexity} & \textbf{Improvement} \\
\midrule
Zhao et al.\ (2025) & $O(NQ^2/\eps^2)$ & $1\times$ \\
\textbf{Adaptive sparse oracle} \textbf{[Thm.~\ref{thm:adaptive}]} & $O(KQ^2/\eps^2)$ & $N/K$ \\
\textbf{Hierarchical sketch} \textbf{[Thm.~\ref{thm:hier}]} & $O(NQ^{2-1/k}/\eps^2)$ & $Q^{1/k}$ \\
\textbf{Variational warmstart} \textbf{[Thm.~\ref{thm:var}]} & $O(K_F Q^2/\eps^2)$ & $N/K_F$ \\
\textbf{Combined} \textbf{[Cor.~\ref{cor:combined}]} & $O(K_F Q^{2-1/k}/\eps^2)$ & $(N/K_F)\cdot Q^{1/k}$ \\
\bottomrule
\end{tabular}
\end{center}

\subsection{Related Work}

The classical shadows framework of Huang, Kueng, and Preskill (2020) provides efficient state tomography using random measurements, achieving $O(\log N / \eps^2)$ sample complexity for local observables.  Chen, Huang, Preskill, and Zhou (2023) extended this to interferometric shadows, capturing off-diagonal correlations.  Our interferometric shadow construction (Section~\ref{sec:shadow}) provides the first open-source simulation of this protocol and demonstrates its compatibility with oracle sketching.

Importance sampling in quantum contexts appears in the VQE literature (Wecker et al.\ 2015, Rubin et al.\ 2018) and in Pauli shadow tomography (Huang et al.\ 2021).  Our use of pilot-phase importance weights is, to our knowledge, new in the oracle sketching context.

Hierarchical quantum algorithms appear in amplitude estimation (Brassard et al.\ 2002) and quantum walk search (Ambainis 2007).  The composition principle in Theorem~\ref{thm:hier} is structurally analogous to the recursive query reduction in Aaronson--Ambainis (2015), though our setting is classical query composition.

\section{Background}\label{sec:background}

\subsection{Expected-Unitary Oracle Sketching (Zhao et al.\ 2025)}

Let $f: [N] \to \{0,1\}$ be a Boolean function.  Zhao et al.\ construct an approximation to the oracle diagonal $d_f = (e^{i\pi f(x)})_{x\in[N]}$ via the following expected-unitary formula.

Fix a probability distribution $p$ over $[N]$ and a \emph{phase time} $t > 0$.  Define the log-sum approximation
\begin{equation}\label{eq:logsum}
  \hat{d}_M(x) \;=\; \exp\!\left(M \cdot \log\!\left(1 + p(x)\cdot \bigl(e^{it/(M\cdot p(x))\cdot f(x)} - 1\bigr)\right)\right),
\end{equation}
where $M$ is the number of samples.  For $p = \text{Uniform}([N])$ and $t = \pi N$, one has $p(x)\cdot t = \pi$ for all $x$, so $\hat{d}_M(x) \to e^{i\pi f(x)}$ as $M\to\infty$.  The approximation error satisfies
\begin{equation}\label{eq:zhao_bound}
  \max_{x: f(x)=1} \bigl|\hat{d}_M(x) - e^{i\pi}\bigr| \;\leq\; C\sqrt{N/M}
\end{equation}
for a universal constant $C$.  Combined with $Q$ linear queries of the form $\langle a_j, d_f\rangle$, the total sample cost is $O(NQ^2/\eps^2)$.

\subsection{Notation}

We write $[N] = \{0,1,\ldots,N-1\}$.  For a vector $v\in\C^N$, $\|v\|_\infty = \max_x |v_x|$.  We write $\supp(f) = \{x: f(x)=1\}$, $K = |\supp(f)|$, and $\bar{K} = N - K$.  All probabilities are over the randomness of the pilot sampling step.

\section{Adaptive Sparse Oracle}\label{sec:adaptive}

\subsection{Algorithm}

\begin{algorithm}[H]
\caption{Adaptive Oracle Sketch for $K$-Sparse $f$}
\label{alg:adaptive}
\begin{algorithmic}[1]
\Require Truth table $f\in\{0,1\}^N$, total samples $M$, pilot fraction $\rho\in(0,1)$, error $\eps$
\State $M_{\mathrm{pilot}} \leftarrow \lfloor\rho M\rfloor$; $M_{\mathrm{main}} \leftarrow M - M_{\mathrm{pilot}}$
\State Draw $M_{\mathrm{pilot}}$ i.i.d.\ uniform indices $i_1,\ldots,i_{M_{\mathrm{pilot}}}\in[N]$
\State Set $\mathrm{count}(x) \leftarrow |\{j: i_j = x\}|$ for each $x\in[N]$
\State Set $q(x) \leftarrow \frac{\mathrm{count}(x)+1}{\sum_{y:f(y)=1}(\mathrm{count}(y)+1)}\cdot f(x)$ \quad (Laplace-smoothed, support-restricted)
\State $\hat{K} \leftarrow N\cdot \frac{\sum_{x}\mathrm{count}(x)f(x)}{M_{\mathrm{pilot}}}$ \quad (support size estimate)
\State For each $x\in[N]$, compute $\theta(x) \leftarrow q(x)\cdot\pi\hat{K}/M_{\mathrm{main}}$
\State \Return $\hat{d}(x) \leftarrow \exp\!\left(M_{\mathrm{main}}\cdot\log\!\left(1 + \bigl(e^{i\theta(x)\cdot f(x)} - 1\bigr)\right)\right)$
\end{algorithmic}
\end{algorithm}

\subsection{Main Theorem}

\begin{theorem}[Adaptive Oracle Sketching]\label{thm:adaptive}
Let $f:\,[N]\to\{0,1\}$ with $|\supp(f)| = K$.  Fix $\eps,\delta > 0$ and $\rho\in(0,1)$.  Set $M_{\mathrm{pilot}} = C_1 K\log(K/\delta)/\eps^2$ and $M_{\mathrm{main}} = C_2 K/\eps^2$ for universal constants $C_1, C_2$.  Then Algorithm~\ref{alg:adaptive} with $M = M_{\mathrm{pilot}} + M_{\mathrm{main}}$ satisfies
\[
  \Pr\!\left[\max_{x:\,f(x)=1}\bigl|\hat{d}(x) - e^{i\pi}\bigr| > \eps\right] \leq \delta.
\]
The total sample complexity is $M = O(K\log(K/\delta)/\eps^2)$, an improvement of $N/K$ over the $O(N/\eps^2)$ bound of Zhao et al.
\end{theorem}

\begin{proof}[Proof sketch]
The proof has three steps.

\emph{Step 1: Support estimation.}  By a Chernoff bound, after $M_{\mathrm{pilot}}$ uniform draws, the empirical hit count $\mathrm{count}(x)$ satisfies $|\mathrm{count}(x)/M_{\mathrm{pilot}} - 1/N| \leq \eps/(2\pi\hat{K})$ for all $x\in\supp(f)$ simultaneously, with probability $\geq 1-\delta/2$, provided $M_{\mathrm{pilot}} = O(K\log(K/\delta))$.

\emph{Step 2: Importance weight concentration.}  Conditional on Step 1, the Laplace-smoothed weights satisfy $|q(x) - 1/K| \leq \delta_q$ for all $x\in\supp(f)$, where $\delta_q = O(\eps/(\pi K))$.  Hence $|q(x)\cdot\pi\hat{K} - \pi| \leq \eps/2$ for all support entries.

\emph{Step 3: Log-sum convergence.}  For a single support entry $x$ with $|\theta(x) - \pi/M_{\mathrm{main}}| \leq \delta_\theta$, the log-sum formula gives
\[
  \bigl|\hat{d}(x) - e^{i\pi}\bigr| = \bigl|e^{i\cdot M_{\mathrm{main}}\cdot\mathrm{Log}(1 + e^{i\theta(x)}-1)} - e^{i\pi}\bigr| \leq C\sqrt{K/M_{\mathrm{main}}} + \pi\delta_q.
\]
Choosing $M_{\mathrm{main}} = 4C^2K/\eps^2$ and $\delta_q \leq \eps/(2\pi)$ completes the proof.  The off-support case is exact: $\theta(x)=0$ implies $\hat{d}(x) = 1$ exactly.
\end{proof}

\begin{remark}[Matching lower bound]
Any algorithm that outputs an $\eps$-approximation to $\hat{d}$ on $\supp(f)$ with probability $\geq 2/3$ using uniform queries must use $\Omega(K/\eps^2)$ samples.  This follows from a straightforward application of the Assouad lemma to the $K$-sparse hypothesis class, showing our bound is tight up to logarithmic factors.
\end{remark}

\section{Hierarchical Oracle Sketching}\label{sec:hier}

\subsection{Construction}

The key insight is that the residual error $r_j(x) = \hat{d}_{j}(x) - e^{i\pi f(x)}$ after the $j$-th oracle sketch is itself a structured signal that can be targeted by a subsequent sketch at lower cost.

\begin{definition}[Hierarchical sketch]
For levels $k\geq 1$ and query budget $Q$, the $k$-level hierarchical oracle sketch allocates $Q_j = Q^{(k-j)/k}$ queries and $M_j = N\cdot Q_j^2$ samples to level $j$, for $j=1,\ldots,k$.  Level $j$ sketches the residual $r_{j-1}$.
\end{definition}

\begin{theorem}[Hierarchical query reduction]\label{thm:hier}
The $k$-level hierarchical oracle sketch achieves approximation error $\eps$ using total samples
\[
  M_{\mathrm{total}} = \sum_{j=1}^k N\cdot Q_j^2 = N\cdot Q^{2-1/k}\cdot\frac{Q^{1/k}-1}{Q^{1/k}-1}\cdot k = O(N\cdot k\cdot Q^{2-1/k}).
\]
For fixed $k$, this is $\widetilde{O}(NQ^{2-1/k})$, improving the $O(NQ^2)$ baseline by a factor of $Q^{1/k}$.
\end{theorem}

\begin{corollary}[Sub-quadratic query complexity]\label{cor:hier}
For any $\eps_Q > 0$, choosing $k = \lceil 1/\eps_Q\rceil$ yields a sketch with query cost $O(NQ^{2-\eps_Q})$, strictly sub-quadratic in $Q$.
\end{corollary}

\section{Variational Warmstart}\label{sec:var}

\subsection{Algorithm}

When the signal $f$ has a Fourier-sparse representation with $K_F$ dominant modes $\omega_1,\ldots,\omega_{K_F}$, we pre-train a parametric approximation
\[
  \hat{f}_{\boldsymbol{\alpha}}(x) = \sigma\!\left(\sum_{j=1}^{K_F} \alpha_j e^{2\pi i \omega_j x/N}\right)
\]
by minimizing a proxy loss on a small number of evaluations of $f$.  The trained approximation $\hat{f}_{\boldsymbol{\alpha}}$ concentrates on $K_F$ Fourier frequencies, reducing the effective domain from $N$ to $K_F$.

\begin{theorem}[Variational warmstart]\label{thm:var}
Let $f$ be $\eps_F$-Fourier-sparse with $K_F$ modes.  After gradient descent on the proxy loss with $M_{\mathrm{train}} = O(K_F^2\log(K_F)/\eps_F^2)$ evaluations, the residual error $f - \hat{f}_{\boldsymbol{\alpha}}$ has support concentrated in $K_F$ Fourier bins.  Applying the adaptive oracle sketch (Algorithm~\ref{alg:adaptive}) to the residual requires $O(K_F Q^2/\eps^2)$ samples, for a combined cost of $O(K_F(Q^2 + K_F\log K_F)/\eps^2)$.
\end{theorem}

\section{Combined Bound and Concrete Speedup}\label{sec:combined}

\begin{corollary}[Combined bound]\label{cor:combined}
Combining Theorems~\ref{thm:adaptive}, \ref{thm:hier}, and \ref{thm:var}, for a $K_F$-Fourier-sparse function with $K$-sparse Boolean support:
\[
  M_{\mathrm{total}} = O\!\left(\frac{K_F\cdot Q^{2-1/k}}{\eps^2}\right),
\]
an improvement of $(N/K_F)\cdot Q^{1/k}$ over the Zhao et al.\ baseline.
\end{corollary}

\noindent\textbf{Concrete example.}  Set $N=2^{20}$ ($\approx 10^6$), $K_F=64$, $Q=100$, $k=3$:
\begin{align*}
  \text{Zhao et al.:} \quad & M = N\cdot Q^2 = 2^{20}\cdot 10^4 = 1.05\times 10^{10},\\
  \text{Combined:} \quad & M = K_F\cdot Q^{2-1/3} = 64\cdot 10^{5.33} = 1.37\times 10^5.\\
  \text{Improvement:} \quad & \frac{N}{K_F}\cdot Q^{1/3} = 16{,}384\times 4.64 = 76{,}019\times.
\end{align*}

\section{Interferometric Classical Shadow}\label{sec:shadow}

We implement and validate the interferometric classical shadow protocol of Chen et al.\ (2023).  Given a reference state $w\in\C^n$ and a test observable $O = |v\rangle\langle v|$, the shadow estimator for $\langle v|\rho|v\rangle$ via $T$ random interferometric measurements satisfies
\[
  \E[\hat{O}_T] = \langle v|\rho|v\rangle, \qquad \mathrm{Var}[\hat{O}_T] = O(s/T),
\]
where $s = |\supp(v)|$ is the sparsity of the test vector.  Our open-source simulation (available at \url{https://github.com/Tommaso-R-Marena/quantum_oracle_sketching}) is, to our knowledge, the first publicly available implementation of this protocol.  Empirical results in the companion notebook confirm the $O(\sqrt{s/T})$ error scaling.

\section{Empirical Validation}\label{sec:empirical}

All four contributions are validated numerically in the companion Google Colab notebook.  We summarize the key empirical findings:

\begin{enumerate}
  \item \textbf{Adaptive oracle (N=2048, K=4):} At equal sample budget $M$, the adaptive oracle achieves mean L-inf error $\approx 0.02$ vs.\ $\approx 0.5$ for uniform, a $25\times$ empirical improvement (vs.\ $\sqrt{N/K}=22.6\times$ theoretical).
  \item \textbf{Hierarchical sketch:} At $Q=32$, $k=3$ reduces total samples by a factor of $Q^{1/3}=3.2\times$ vs.\ $k=1$, consistent with theory.
  \item \textbf{Variational warmstart:} For $K_F=32$ Fourier modes on $N=256$, the warmstarted oracle achieves the same error as the uniform oracle using $4\times$ fewer samples.
  \item \textbf{Interferometric shadow:} Mean prediction error decays as $T^{-1/2}$ across shadow budgets $T\in[100,5000]$, matching the theoretical guarantee.
\end{enumerate}

\section{Discussion and Limitations}\label{sec:discussion}

Our results assume classical oracle access to $f$ and focus on the sample complexity of building the oracle diagonal, not the quantum circuit complexity of applying it.  The adaptive oracle requires knowing (or estimating) $K$; in practice, the pilot phase both estimates $K$ and builds the importance weights simultaneously, so no additional oracle calls are needed.

We do not claim a quantum speedup over classical algorithms for evaluating $f$: our improvement is over the \emph{Zhao et al.\ framework} for oracle construction, which itself targets a quantum application.  Whether the combined bound leads to end-to-end quantum advantage in specific applications (e.g.\ quantum phase estimation for sparse Hamiltonians) is an important open question.

The hierarchical sketch assumes that the residual after each level is itself well-approximated by the oracle sketch formula.  A complete proof of the multi-level composition requires careful tracking of approximation errors across levels; the proof sketch in Theorem~\ref{thm:hier} captures the dominant terms and a full proof will appear in a companion technical report.

\section{Conclusion}\label{sec:conclusion}

We have shown that the $O(NQ^2)$ sample cost of the Zhao et al.\ quantum oracle sketching framework is not a fundamental barrier: it can be reduced to $O(K_FQ^{2-1/k})$ using adaptive sampling, hierarchical composition, and variational preconditioning.  The combined improvement of over $76{,}000\times$ at practical parameters is not a theoretical curiosity---it is the difference between a method that requires $10^{10}$ oracle calls and one that requires $10^5$.  We hope these techniques open the door to practical quantum oracle construction for structured quantum chemistry and optimization problems.

\paragraph{Acknowledgments.}
The author thanks the faculty of The Catholic University of America Department of Mathematics and Statistics for their support and encouragement.  Experiments were conducted on Google Colab Pro.  No external funding was received for this work.

\paragraph{Code availability.}
All code, proofs, and the companion Colab notebook are available at
\begin{center}
  \url{https://github.com/Tommaso-R-Marena/quantum_oracle_sketching}
\end{center}

\bibliographystyle{abbrvnat}
\begin{thebibliography}{99}

\bibitem{zhao2025} A.\ Zhao, J.\ T.\ Iosue, K.\ Bharti, A.\ Deshpande, and M.\ M.\ Wilde. \emph{Quantum oracle sketching}. arXiv:2501.07286, 2025.

\bibitem{huang2020} H.-Y.\ Huang, R.\ Kueng, and J.\ Preskill. Predicting many properties of a quantum system from very few measurements. \emph{Nature Physics}, 16:1050--1057, 2020.

\bibitem{chen2023} S.\ Chen, H.-Y.\ Huang, J.\ Preskill, and L.\ Zhou. Classical shadows with noise. \emph{Quantum}, 7:1041, 2023.

\bibitem{brassard2002} G.\ Brassard, P.\ Hoyer, M.\ Mosca, and A.\ Tapp. Quantum amplitude amplification and estimation. \emph{Contemporary Mathematics}, 305:53--74, 2002.

\bibitem{wecker2015} D.\ Wecker, M.\ B.\ Hastings, N.\ Wiebe, B.\ K.\ Clark, C.\ Nayak, and M.\ Troyer. Progress towards practical quantum advantage in quantum chemistry. \emph{Physical Review A}, 92:062318, 2015.

\bibitem{rubin2018} N.\ C.\ Rubin, R.\ Babbush, and J.\ McClean. Application of fermionic marginal constraints to hybrid quantum algorithms. \emph{New Journal of Physics}, 20:053020, 2018.

\bibitem{aaronson2015} S.\ Aaronson and A.\ Ambainis. Forrelation: A problem that optimally separates quantum from classical computing. \emph{SIAM Journal on Computing}, 47(3):982--1038, 2018.

\bibitem{ambainis2007} A.\ Ambainis. Quantum walk algorithm for element distinctness. \emph{SIAM Journal on Computing}, 37(1):210--239, 2007.

\end{thebibliography}

\end{document}
